% Figure 43.2 : un rôle, une liaison, des sujets ; trois combinaisons (chapitre 43).
\documentclass[tikz,border=4pt]{standalone}
\usepackage{quai}
\begin{document}
\begin{tikzpicture}[x=1cm,y=1cm,
    b/.style={qbox, font=\scriptsize, align=left, inner sep=3pt, text width=3.6cm, minimum height=1cm},
    lien/.style={qlbl, font=\scriptsize, fill=QPaper, inner sep=1pt, align=center}]

  \node[qtitre] at (-1.9,1.2) {sujets};
  \node[qtitre] at (3,1.2) {liaison};
  \node[qtitre] at (8.2,1.2) {rôle};

  \foreach \y/\c/\s/\l/\r/\p in {
      0/QVert/{groupe {\ttfamily equipe-colis}}/{{\ttfamily RoleBinding astreinte}\\namespace {\ttfamily colis}}/{{\ttfamily Role astreinte}\\défini dans {\ttfamily colis}}/{valable dans {\ttfamily colis}},
      -2/QBleu/{utilisatrice {\ttfamily carla}}/{{\ttfamily RoleBinding cours-view-carla}\\namespace {\ttfamily colis}}/{{\ttfamily ClusterRole view}\\défini pour tout le cluster}/{valable dans {\ttfamily colis} seulement},
      -4/QOcre/{ServiceAccount\\{\ttfamily kube-system/default}}/{{\ttfamily ClusterRoleBinding}\\{\ttfamily minikube-rbac}}/{{\ttfamily ClusterRole cluster-admin}}/{valable partout}} {
    \node[b, draw=\c, fill=\c Bg] (s) at (-0.6,\y) {\s};
    \node[b, draw=\c] (l) at (4.6,\y) {\l};
    \node[b, draw=\c] (r) at (9.8,\y) {\r};
    \draw[qarr] (l.west) -- node[lien, above] {subjects} (s.east);
    \draw[qarr] (l.east) -- node[lien, above] {roleRef} (r.west);
    \node[qlbl, font=\scriptsize, text=\c, anchor=west] at (11.9,\y) {\p};
  }
\end{tikzpicture}
\end{document}
